\documentclass[../../main.tex]{subfiles}
\begin{document}

{\bf [解答]}

条件(1)から$f'(0), g'(0)$が存在し，条件(3)から
\begin{align}
  f'(0)g'(0) = -1 \label{eq:1}
\end{align}
である．

次に条件(4)から
\begin{align}
  &
  \begin{cases}
    f'(0)=\lim_{x\to0}\dfrac{a+bf(x)/x}{c+f(x)/x} \\
    g'(0)=\lim_{x\to0}\dfrac{a+bg(x)/x}{c+f(x)/x}
  \end{cases} \\
  &
  \begin{cases}
    f'(0)=\dfrac{a+bf'(0)}{c+f'(0)} \\
    g'(0)=\dfrac{a+bg'(0)}{c+g'(0)}
  \end{cases}
\end{align}
である．よって$f'(0),g'(0)$は
\begin{align}
  x = \dfrac{a+bx}{c+x}
  \quad \therefore
  x^2 + (c-b)x -a = 0 \label{eq:2}
\end{align}
の解である．\cref{eq:1}より$f'(0)\neq g'(0)$であることと，
これは二次方程式でたかだか二つしか解がないことから，
$f'(0), g'(0)$は$2$異実解を構成する．
よって解と係数の関係と\cref{eq:1}から
\begin{align}
  -a = f'(0)g'(0) = -1
  \quad\therefore
  a = 1
\end{align}
である．

\newpage
\end{document}
